% Section 11: numerical safeguards: what went wrong and how each fix works.
\section{Numerical safeguards: what went wrong and what fixed it}\label{sec:safeguards}

Most of the effort in a simplex or interior-point code goes into its failure modes, not into the
textbook iteration. This section collects the failures we met on the benchmark sets and explains
the mechanism of each fix. The first sweeps (\file{netlib\_simplex\_v1.csv}, \file{netlib\_ipm\_v1.csv})
record the state before the fixes. \Cref{tab:fixes-spx,tab:fixes-ipm} list every problem whose result
changed between versions.

\subsection{Dual simplex}\label{sec:simplex-safeguards}

\begin{table}[tbp]
\centering\small
\caption[Dual simplex: problems changed by the fixes]{Dual simplex, v1 against v2: every Netlib
problem not solved by both versions. v1 solved \nFixedSpxVOne{} of \nFixedSpxVOneTot{}; the fixes of
\cref{sec:simplex-safeguards} solved \nFixedSpx{} more. On the problems both versions finish, v2 is
\nFixedSpxSpeedup{} times faster in shifted geometric mean, measured relative to HiGHS in each sweep.}
\label{tab:fixes-spx}
\begin{tabular}{@{}l lrr lrr@{}}
\toprule
 & \multicolumn{3}{c}{v1} & \multicolumn{3}{c}{v2}\\
\cmidrule(lr){2-4}\cmidrule(l){5-7}
Problem & status & rel.\ err & time (s) & status & rel.\ err & time (s)\\
\midrule
\rows{data/fixes_simplex.tex}
\bottomrule
\end{tabular}
\end{table}

\paragraph{False ``infeasible'' from the ratio test (\code{perold}, \code{pilotnov})}
The v1 sweep declared two feasible problems infeasible within a second. A dual simplex concludes
primal infeasibility when the ratio test finds no entering candidate: the dual is unbounded along the
pivot row, which is a Farkas certificate. That conclusion is only as good as the reduced costs and the
pivot row it rests on. The v1 ratio test did not use the textbook Harris bound. A Harris pass must
bound the step by the \emph{smallest relaxed} ratio $\min_j(|d_j|+\epsilon_d)/|\alpha_{rj}|$ over the
candidates that remain after the bound flips. Only then does a step to any candidate with $t_j\le\Theta$
leave every reduced cost within $\epsilon_d$ of its feasible sign. A step past that bound makes some
$d_j$ infeasible by more than the tolerance. Later ratio tests then exclude that $j$ or treat it
wrongly, and an empty candidate set, and so a false certificate, follows. The fixed kernel
(\cref{alg:ratio}) computes the bound exactly as above. It also makes the pivot tolerance
\emph{relative} to the largest $|\alpha_{rj}|$ in the row, because the \code{pilot} models have rows
whose entries span many orders of magnitude, where an absolute $10^{-7}$ admits pivots that are noise.
Two further guards keep a certificate honest. First, before reporting infeasibility the code
refactorises, recomputes $x_{\mathcal B},y,d$ from scratch, repairs the dual feasibility, and repeats the
iteration; only a second empty ratio test on a fresh factorisation counts. Second, the pivot seen
from the row ($\alpha_{rq}$) and from the column ($\alpha_{q,r}$) must agree to $10^{-6}$ relative, or
the factorisation is renewed.

The source of v1 is not in the repository: the first commit already contains the fixed kernel. This
account is therefore of the fixed code and of the failure recorded in the v1 sweep.

\paragraph{Stalling and cycling in phase 1 (\code{cycle}, \code{pilot87}, \code{tuff})}
Three problems ran into the \tlNetlib\,s limit in v1. Degenerate problems give many zero
reduced costs, so many ratio-test breakpoints tie at zero, and the dual simplex can pivot through a
sequence of degenerate bases without improving the dual objective, possibly returning to a basis
already visited. v1 perturbed the costs in phase~2 only, while phase~1 on Fourer's auxiliary problem
\eqref{eq:aux} is at least as degenerate: every bound there is $0$ or $\pm1$. The fix applies the
random perturbation \eqref{eq:perturb} in phase~1 as well, which breaks the ties. \code{cycle}
went from a time-out to under a second (\cref{tab:fixes-spx}).

\paragraph{Wrong objective on \code{pilot}: the dual tolerance of the cleanup}
v1 returned \code{optimal} on \code{pilot} with a relative objective error far above $10^{-6}$
(\cref{tab:fixes-spx}). The cause is the size of the variable ranges. After the perturbation is removed, reduced costs with the
wrong sign but below the dual tolerance $\epsilon_d=10^{-7}$ are accepted as optimal. The objective
error they leave is bounded by
\begin{equation}\label{eq:dualerr}
  |f(x)-f^\star| \;\lesssim\; \sum_{j\in\mathcal N} |d_j^{\text{infeas}}|\,(u_j-\ell_j),
\end{equation}
and on \code{pilot} a reduced cost of $10^{-8}$ on a variable with a scaled range of $10^3$ already
shifts the objective by $10^{-5}$. The fix polishes with the primal simplex to $\epsilon_d=10^{-9}$ in
the cleanup phase only (phase~1 and phase~2 keep $10^{-7}$, which is enough to steer the search).
v2 solves \code{pilot} to the reference objective (\cref{tab:fixes-spx}).

\paragraph{Singular bases}
Basis repair (\cref{sec:lu}) replaces a dependent basic column by a logical instead of failing; the
number of repairs is reported with each result. Refactorisation every 100 updates, or when the eta
file grows beyond eight times the LU size, bounds the growth of rounding error in the product form.

\paragraph{Exact duals at return}
The iteration updates $d$ incrementally and $y$ only at refactorisation, and cost shifts change $c$
during the loop. The returned $y$ is therefore recomputed as $B\mT c_{\mathcal B}$ for the true costs.
This matters for postsolve (\cref{sec:presolve-bug}) and for anyone who uses the duals as prices, as
refinery planners do.

\subsection{Interior-point method}

\begin{table}[tbp]
\centering\small
\caption[Interior point: problems changed by the fixes]{Interior point, v1 against v2: every Netlib
problem not solved by both versions. v1 solved \nFixedIpmVOne{} of \nFixedIpmVOneTot{}; v2 solves
\nFixedIpm{} more. On problems both versions finish, v2 is \nFixedIpmSpeedup{} times faster in shifted
geometric mean, measured relative to HiGHS in each sweep.}
\label{tab:fixes-ipm}
\begin{tabular}{@{}l lrr lrr@{}}
\toprule
 & \multicolumn{3}{c}{v1} & \multicolumn{3}{c}{v2}\\
\cmidrule(lr){2-4}\cmidrule(l){5-7}
Problem & status & rel.\ err & time (s) & status & rel.\ err & time (s)\\
\midrule
\rows{data/fixes_ipm.tex}
\bottomrule
\end{tabular}
\end{table}

\paragraph{Stall detection}
In v1 several problems (\code{degen3}, \code{greenbeb}, \code{pilot}, \code{pilot87}, \code{dfl001})
ran to the iteration limit, some of them with objectives already correct to $10^{-8}$. Near the
solution $\Theta$ spans many orders of magnitude and the normal equations lose the digits
needed to reach $10^{-8}$. Further iterations only drive $\mu$ towards zero and $z_j/x_j$ towards
overflow. The v2 rule of \cref{sec:ipm} detects this. When the error has not halved in four iterations
(or $\mu$ has collapsed), a point that meets $10^{-6}$ is accepted, and after 30 such iterations the
method gives up with \code{numerical\_error} instead of wandering. Residuals are measured as relative
$\infty$-norms \eqref{eq:ipm-term}, the convention of HiGHS/IPX, so the tolerance means the same on
every problem. The remaining failures (\ipmFailNames) are the problems where even $10^{-6}$ is out of
reach of the normal equations without presolve and crossover (\cref{sec:results}).

\paragraph{Cholesky pivot guard}
Dependent rows make $M$ singular at the limit. A pivot $\le10^{-30}$ is replaced by $10^{128}$
(\cref{sec:chol}), which zeroes that component of $\Delta y$ instead of producing \code{inf}. The dual
regularisation $\delta=10^{-10}$ and the two refinement steps keep the resulting error small. The
numerical errors on \code{grow15} and \code{grow22} in v1 no longer occur.

\subsection{QP interior point: the \texorpdfstring{$LDL\T$}{LDLT} pivot floor}
The augmented matrix is quasi-definite only in exact arithmetic with $\rho,\delta>0$. When $D$
spans many orders of magnitude, cancellation can still produce a pivot of the wrong sign or of
magnitude $10^{-20}$, and dividing by it destroys the factorisation. The kernel replaces any pivot with
$\sigma_kD_{kk}<\delta_{\min}$ by $\sigma_k\delta_{\min}$ (\cref{sec:ldl}), so the inertia is always
right. Without the floor a single bad pivot produces \code{nan} in the direction and ends the solve.
Refinement against the true matrix recovers the accuracy \eqref{eq:refine}. The floor alone is not
enough, though. Flooring many pivots changes the matrix so much that the direction can be finite and
still wrong (\code{QRECIPE}). The method therefore also checks the refined residual of every direction,
and if the check fails it boosts the regularisation by $10^2$ and refactorises (\cref{sec:qp}).

\subsection{Branch-and-bound: heap tie-breaking}\label{sec:heap}
When the objective is integral, node bounds are rounded up, so many open nodes share exactly the
same key. A heap keyed by $(\underline z,\text{counter})$ breaks those ties by insertion order: it
takes the \emph{oldest}, hence shallowest, node among equals. The search then proceeds breadth-first
across a level of equal bounds, the open list grows, and incumbents that a deeper node would give
arrive late. Adding $-\text{depth}$ as the second key \eqref{eq:heapkey} makes ties go to the deepest
node. This changes nothing about which bound is processed next, so the best-bound guarantee is
unchanged. It only makes the search finish subtrees it has already started.

\subsection{Presolve: dual postsolve order}
See \cref{sec:presolve-bug}: singleton-row duals depended on the order in which reductions were
undone, and recording the column's cost and live rows at removal time fixed it.

\begin{lessonbox}[What these failures have in common]
The worst of these bugs produced confident wrong answers (``infeasible'', ``optimal'' with the wrong
objective, duals with the wrong sign) rather than crashes or error messages. Each was found because the benchmark
harness checks every result against an independent solver in a separate process, and because the dual
check tests the optimality conditions on the original model. \Cref{sec:verification} describes these
checks.
\end{lessonbox}
